\subsectionaddtolist{ب}
برای زبان $L_2 = \{ a^n b^m c^k \mid n \le m \le k \}$، یک ماشین تورینگ به صورت زیر تعریف می‌کنیم:

\begin{enumerate}
    \item {بررسی صحت فرمت ورودی:} ابتدا ماشین یک بار نوار را از چپ به راست می‌خواند تا مطمئن شود رشته ورودی دقیقا به فرم تعدادی $a$، سپس تعدادی $b$ و در نهایت تعدادی $c$ است. در صورت مشاهده هرگونه ترتیب غلط (مثلا وجود $a$ بعد از $b$)، ماشین ورودی را رد می‌کند. سپس هد به ابتدای نوار برمی‌گردد.
    
    \item {بررسی شرط $n \le m$:} ماشین از ابتدای نوار شروع می‌کند:
    \begin{itemize}
        \item اولین نماد $a$ را پیدا کرده و آن را با یک نماد جدید (مثلا $X$) جایگزین می‌کند.
        \item به سمت راست حرکت می‌کند تا به اولین نماد $b$ علامت‌نخورده برسد. آن را با یک نماد جدید (مثلا $Y$) جایگزین می‌کند.
        \item سپس به سمت چپ و به ابتدای نوار برمی‌گردد و این روند را تکرار می‌کند.
        \item اگر ماشین نماد $a$ را خواند ولی در مسیر راست هیچ نماد $b$ علامت‌نخورده‌ای پیدا نکرد (به $c$ یا انتهای نوار رسید)، به این معنی است که $n > m$ بوده و ورودی را رد می‌کند.
        \item اگر ماشین تمام $a$ها را با موفقیت با $b$ها متناظر کرد و دیگر هیچ $a$ای روی نوار باقی نماند، هد را مجددا به ابتدای نوار می‌برد و وارد مرحله بعد می‌شود.
    \end{itemize}

    \item {بررسی شرط $m \le k$:} اکنون تمام نمادهای $b$ (چه آن‌هایی که دست‌نخورده‌اند و چه آن‌هایی که به $Y$ تبدیل شده‌اند) باید با نمادهای $c$ متناظر شوند.
    \begin{itemize}
        \item ماشین نوار را جستجو می‌کند تا اولین نماد $b$ یا $Y$ را پیدا کند و آن را با یک نماد جدید (مثلا $Z$) جایگزین کند.
        \item به سمت راست حرکت می‌کند تا به اولین نماد $c$ علامت‌نخورده برسد. آن را با یک نماد جدید (مثلا $W$) جایگزین می‌کند.
        \item به سمت چپ برمی‌گردد و این کار را تکرار می‌کند.
        \item اگر ماشین یک نماد $b$ یا $Y$ را خواند ولی در ادامه هیچ نماد $c$ای برای خط زدن پیدا نکرد، به این معنی است که $m > k$ بوده و ورودی را رد می‌کند.
        \item اگر تمام نمادهای $b$ و $Y$ با موفقیت با $c$ها متناظر شدند، ماشین متوقف شده و رشته را می‌پذیرد.
    \end{itemize}
\end{enumerate}

\begin{tikzpicture}[
    shorten >=1pt, 
    auto, 
    thick, 
    >=Stealth,
    every edge node/.style={align=center, font=\small, fill=white, inner sep=2pt}
]

    % نقطه شروع
    \node[state, initial, initial text=] (q0) at (-4, 0) {$q_0$};
    
    % فاز اول (فشرده‌تر شده)
    \node[state] (q1) at (0, 3) {$q_1$};
    \node[state] (q2) at (-4, 3) {$q_2$};
    
    % فاز دوم (کشیده شده به سمت چپ)
    \node[state] (q4) at (2, 0) {$q_4$};
    \node[state] (q5) at (7.5, 0) {$q_5$};
    \node[state] (q3) at (4.75, -2.5) {$q_3$};
    
    % فاز تایید نهایی و پذیرش 
    \node[state] (q6) at (4.75, -10) {$q_6$};
    \node[state, accepting] (qacc) at (8.5, -5) {$q_{acc}$};

    % یال‌های فاز اول
    \path[->]
        (q0) edge node[left] {$a \to X, R$} (q1)
        (q1) edge [loop right] node {$a \to a, R$ \\ $Y \to Y, R$} (q1)
        (q1) edge node[above] {$b \to Y, L$} (q2)
        (q2) edge [loop left] node {$a \to a, L$ \\ $Y \to Y, L$} (q2)
        (q2) edge node[below left] {$X \to X, R$} (q0);

    % انتقال از فاز اول به فاز دوم
    \path[->]
        (q0) edge node[above] {$Y \to Z, R$ \\ $b \to Z, R$} (q4);

    % یال‌های فاز دوم
    \path[->]
        (q4) edge [loop above] node {$Y \to Y, R$ \\ $b \to b, R$ \\ $W \to W, R$} (q4)
        (q4) edge node[above] {$c \to W, L$} (q5)
        (q5) edge [loop above] node {$Y \to Y, L$ \\ $b \to b, L$ \\ $W \to W, L$} (q5)
        (q5) edge node[below right] {$Z \to Z, R$} (q3)
        (q3) edge node[below left] {$Y \to Z, R$ \\ $b \to Z, R$} (q4);

    % یال‌های تایید نهایی و پایان
    \path[->]
        (q0) edge node[below left, pos=0.6] {$c \to c, R$} (q6)
        
        % قوس تنظیم شده برای جلوگیری از برخورد با q6
        (q0) edge [out=-90, in=210, looseness=1.3] node[below left] {$\sqcup \to \sqcup, R$} (qacc)
        
        (q3) edge node[left] {$c \to c, R$ \\ $W \to W, R$} (q6)
        (q3) edge node[above right] {$\sqcup \to \sqcup, R$} (qacc)
        
        (q6) edge [loop below] node {$c \to c, R$ \\ $W \to W, R$} (q6)
        (q6) edge node[above] {$\sqcup \to \sqcup, R$} (qacc);

\end{tikzpicture}


\subsectionaddtolist{د}
برای زبان $L_4 = \{ \langle x, y, z \rangle \mid x = y + z, \Sigma = \{1\} \}$، با توجه به اینکه الفبا فقط شامل نماد ۱ است، ورودی‌ها به صورت سیستم اعداد یکانی روی نوار قرار دارند. فرض متعارف این است که این سه ورودی با یک نماد جداکننده مانند $\#$ از هم جدا شده‌اند، یعنی نوار در ابتدا حاوی رشته $1^x \# 1^y \# 1^z$ است. 

تعریف ماشین تورینگ:

\begin{enumerate}
    \item {بررسی صحت فرمت:} ماشین ابتدا نوار را پیمایش می‌کند تا مطمئن شود ورودی دقیقا شامل سه بخش از نمادهای ۱ است که با دقیقا دو نماد $\#$ از هم تفکیک شده‌اند. سپس هد به ابتدای نوار برمی‌گردد.

    \item {متناظر کردن $y$ با $x$:} 
    \begin{itemize}
        \item ماشین به سمت راست حرکت می‌کند تا از اولین $\#$ عبور کرده و به اولین نماد ۱ در بلوک $y$ برسد. آن را خط می‌زند (مثلا با نماد $X$ جایگزین می‌کند).
        \item سپس به سمت چپ برمی‌گردد تا به اولین نماد ۱ در بلوک $x$ برسد و آن را نیز خط می‌زند ($X$ می‌کند). 
        \item اگر در زمان تلاش برای خط زدن از بلوک $x$، متوجه شد که هیچ نماد ۱ در این بلوک باقی نمانده است، یعنی $x$ کوچک‌تر از مجموع است و ورودی را رد می‌کند.
        \item این روند رفت و برگشتی تا زمانی ادامه می‌یابد که تمام نمادهای ۱ در بلوک $y$ خط بخورند.
    \end{itemize}

    \item {متناظر کردن $z$ با $x$:}
    \begin{itemize}
        \item ماشین دقیقا عملیات مشابه مرحله قبل را برای بلوک $z$ انجام می‌دهد. یعنی به سمت راست می‌رود تا از دومین $\#$ عبور کند، اولین نماد ۱ در بلوک $z$ را پیدا کرده و خط می‌زند، سپس به بلوک $x$ برگشته و یک نماد ۱ دیگر را خط می‌زند.
        \item باز هم اگر نماد ۱ در بلوک $x$ برای خط زدن پیدا نشد، ورودی رد می‌شود.
        \item این کار تا خط خوردن تمامی ۱های بلوک $z$ تکرار می‌شود.
    \end{itemize}

    \item {تصمیم‌گیری نهایی:} 
    پس از اینکه تمام نمادهای ۱ در بلوک‌های $y$ و $z$ به پایان رسیدند، ماشین بلوک $x$ را بررسی می‌کند.
    \begin{itemize}
        \item اگر در بلوک $x$ هیچ نماد ۱ باقی نمانده بود (همه $X$ شده بودند)، به این معنی است که تعداد ۱های بلوک اول دقیقا برابر با مجموع دو بلوک دیگر بوده و تساوی $x = y + z$ برقرار است. پس ماشین ورودی را می‌پذیرد.
        \item اگر در بلوک $x$ هنوز نماد ۱ خط‌نخورده وجود داشت، یعنی $x > y + z$ بوده و در این حالت ماشین ورودی را رد می‌کند.
    \end{itemize}
\end{enumerate}


\begin{tikzpicture}[
    shorten >=1pt, 
    node distance=2.5cm and 3cm, 
    on grid, 
    auto, 
    thick, 
    >=Stealth,
    every edge node/.style={align=center, font=\small}
]

    % States
    \node[state, initial, initial text=] (q0) {$q_0$};
    \node[state] (q1) [right=3.5cm of q0] {$q_1$};
    \node[state] (q2) [right=3.5cm of q1] {$q_2$};
    \node[state] (q4) [right=3.5cm of q2] {$q_4$};

    \node[state] (q3) [below=5cm of q2] {$q_3$}; % Return from y
    \node[state] (q5) [below=3cm of q4] {$q_5$}; % Return from z

    \node[state] (q6) [below=4cm of q0] {$q_6$};
    \node[state, accepting] (qacc) [left=4cm of q6] {$q_{acc}$};

    % Crossing out x
    \path[->]
        (q0) edge node {$1 \to X, R$} (q1)
        (q0) edge node [left] {$\# \to \#, R$} (q6);

    % Finding 1 in y block
    \path[->]
        (q1) edge [loop above] node {$1 \to 1, R$} (q1)
        (q1) edge node {$\# \to \#, R$} (q2)
        (q2) edge [loop above] node {$Y \to Y, R$} (q2)
        (q2) edge node {$1 \to Y, L$} (q3)
        (q2) edge node {$\# \to \#, R$} (q4); % y is empty, go to z

    % Returning from y to x
    \path[->]
        (q3) edge [loop below] node {$Y \to Y, L$ \\ $\# \to \#, L$ \\ $1 \to 1, L$} (q3)
        (q3) edge node {$X \to X, R$} (q0);

    % Finding 1 in z block (when y is exhausted)
    \path[->]
        (q4) edge [loop above] node {$Z \to Z, R$} (q4)
        (q4) edge node {$1 \to Z, L$} (q5);

    % Returning from z to x
    \path[->]
        (q5) edge [loop below] node {$Z \to Z, L$ \\ $\# \to \#, L$ \\ $Y \to Y, L$ \\ $1 \to 1, L$} (q5)
        (q5) edge node {$X \to X, R$} (q0);

    % Final Verification (x is exhausted, checking y and z)
    \path[->]
        (q6) edge [loop below] node {$Y \to Y, R$ \\ $Z \to Z, R$ \\ $\# \to \#, R$} (q6)
        (q6) edge node [left] {$\sqcup \to \sqcup, R$} (qacc);

\end{tikzpicture}